\section{Protections and the length of high-type paths}\label{sec:weights}

The structural exclusions give nested simple groups whose actions remain
transitive or absolutely irreducible in every earlier ancestor. We now
need two further conclusions. First, within each consecutive
same-characteristic classical run, the chart construction can select
descendants whose actions on each ancestor's natural module admit no
nontrivial factorization into individually invariant tensor factors
over the algebraic closure, and that are fixed by no prime compositional
root of its defining field map. Second, any path with these protections
has absolutely bounded length. The selection uses the affine geometry;
the length bound uses only the group properties stated below, so it can
also be applied to the later reverse and forward constructions.

All high-degree and high-rank thresholds in this section are sufficiently
large absolute constants. Bounded types have already been treated in
\Cref{sec:bounded}.

\begin{definition}[Prepared paths and ancestor properties]\label{w:prepared-path}
A \emph{prepared path} is a strict descending chain of finite groups
$K_W$ containing a common subgroup $H$, together with actual strictly
nested nonabelian simple groups $B_W\trianglelefteq K_W$ such that
\[
 HB_W=K_W,\qquad C_{K_W}(B_W)\le H,
\]
where each centralizer is soluble and the quotient by it is almost
simple with socle $B_W$. Fix the usual natural realizations of the
alternating and classical socles. For every ancestor--descendant pair
require the following properties.
\begin{enumerate}
\item In an alternating ancestor, the descendant socle is transitive on the natural set.
\item In a high classical ancestor, the descendant socle is projectively absolutely irreducible on the natural module. Its label normalizes no proper ground-field tensor system of the endomorphism algebra of the kind excluded in \Cref{sec:classical}; in isometry types the factors are individually selfadjoint.
\end{enumerate}
\end{definition}

The pruning construction of \Cref{sec:patterns} supplies these properties
by \Cref{c:classical-conclusion,c:alternating-conclusion}.
\Cref{c:ancestor-socles} makes the simple groups actual and compatible
in all ancestor quotients, whose centralizer kernels lie in $H$.
Consequently a full stabilizer containing a tested label pulls back to
a label of the exact original interval. In the selection argument the
same properties hold for all comparable candidate subflats of a
successful pruned flat before the next node is chosen: if
$W\subset L\subset U$, then $B_U\le B_L\le B_W$.

\subsection{Classical conventions and changes of type}

\begin{convention}[High classical groups]\label{w:classical-conventions}
We use the standard structure theory of finite simple classical groups, simply connected coverings, Schur multipliers and automorphisms. In sufficiently high rank there are no small covering exceptions: the universal central extension is the usual fixed group \(\widehat S=\widehat G^\Phi\) of simply connected type. Write \(k=\overline{\mathbf F}_p\) in characteristic \(p\) defining the group. We also use the adjoint algebraic group \(G\) over \(k\), with \(S=(G^\Phi)'\), and regard it on points as a projective group on a natural \(v\)-space. Here \(\Phi\) is standard of parameter \(P=p^f\) (split Frobenius composed possibly with diagram/inner automorphisms); ground field \(F\) on the finite natural module is of size \(P\), or \(P^2\) for unitary type. We use types A (linear and unitary), C, odd orthogonal B only in odd characteristic, and even orthogonal D; by the bilateral (bilinear) form types we mean B,C,D, not the unitary case. The usual automorphisms of \(S\) lift to point maps in the group on adjoint \(G(k)\) generated by inner and diagram automorphisms and integral powers of a split Frobenius \({\rm Fr}_p\) (including inverse on points). Modulo inner this group is abelian, with at most two diagram choices and field exponent called height. A positive-height element is an ordinary Steinberg map (inner conjugate to standard form by Lang--Steinberg). We shall use that in this adjoint group the order of the derived fixed simple group at positive height \(f_2\) exceeds the order of the whole fixed group at \(0<f_1<f_2\), uniformly in high rank. Indeed the classical order formulae give order \((p^f)^{\dim G}\) times factors bounded above and away from zero (products of \(1\pm(p^f)^{-j}\)); the derived index costs at most a polynomial in the rank. When finite types \(B_e,C_e\) coincide in characteristic two we use the symplectic realization (type C); the odd-dimensional \(B_e\) convention in this section is only for odd characteristic. In particular all Steinberg maps needed \emph{in this high classical section} are ordinary graph/field ones, not the special square-root-field maps. We use the standard high-rank multiplier and automorphism results in \citet{GLS1994,KleidmanLiebeck1990}.
\end{convention}

We use the standard highest weight theory of simply connected semisimple algebraic groups, including Steinberg's Tensor Product Theorem and his parametrization of the irreducibles of \(\widehat S\) in defining characteristic as restrictions of \(L_{\widehat G}(\lambda)\) for \(P\)-restricted dominant weights (including standard diagram twists). Projective representations of \(S\) lift to ordinary ones on its universal cover, uniquely once the projective matrices are fixed (perfection). 

More explicitly, we use the universal-central-extension and automorphism description in the high-rank range of the standard multiplier/automorphism theorems for finite groups of Lie type. Thus lifting a projective representation means pulling back the scalar central extension and then using universality on \(\widehat S\). The Steinberg theorems we use here are the ordinary \(p\)-adic tensor product for rational simples and the defining-characteristic Restriction Theorem on the full simply connected fixed group: highest weights with fundamental coefficients \(0,\ldots,P-1\) give irreducibles, pairwise inequivalent, and give all of them, also for these standard twisted forms. With pinned map a diagram twist times split \({\rm Fr}_p^f\), a further field twist rotates digit positions on restriction modulo \(f\), with wrapping applying a diagram (since the defining map is the identity on the finite points). One obtains this by twisting the tensor product digit factors separately. Equivalence of linear lifts up to a projective intertwiner on the perfect cover already implies ordinary equivalence (the scalar discrepancy would be a character). These are the ordinary tensor-product and finite-restriction theorems, including ordinary diagram twists, of \citet[Theorems 1.1, 7.4 and 9.3]{Steinberg1963}; see also \citet{Jantzen2003}.

\begin{lemma}[A lower bound for projective degree]\label{w:min-degree}
There is an absolute constant $c_0>0$ such that every nontrivial projective representation in defining characteristic of a sufficiently high-rank simple classical group with natural dimension $v$ has dimension at least $c_0v$.
\end{lemma}
\begin{proof}
Any nontrivial irreducible projective representation of \(S\) in characteristic \(p\) has dimension at least a positive absolute constant times \(v\) in high rank: its highest-weight Weyl orbit already has at least that size (permutations of a nonconstant weight word in type A; signed permutations, possibly with even sign changes, in B,C,D). This bound suffices here and also bounds any nontrivial projective representation, using perfection if the lift has only trivial composition factors. In the last case a lift with all composition factors trivial has unipotent, hence soluble, image. A perfect group has no nontrivial soluble image, so a nontrivial representation must have a nontrivial composition factor.
\end{proof}

\begin{lemma}[Alternating subsequences]\label{w:alternating-switch}
A prepared path contains at most two alternating socles of sufficiently large degree.
\end{lemma}
\begin{proof}
Any alternating subsequence of the path has bounded length. Suppose there are three socles \(A_n>A_m>A_e\) as subgroups (\(n>m>e\)), with the transitivity for each ancestor pair. For large \(e\), \(m\ge\binom e2\). In fact an intransitive point stabilizer in natural \(A_e\)-coordinates giving index \(>e\) has index at least \(\binom e2\) (if within a point group \(A_{e-1}\) properly, use simplicity for minimum proper index at least \(e-1\)). Proper transitive subgroups of \(A_e\) have index at least exponential in \(e\) for large \(e\), by Bochert's Theorem in the primitive case, and the block partition count in the imprimitive case. For instance with \(b\) blocks of size \(a\), both at least two, pair the blocks disjointly and interchange corresponding positions arbitrarily in each pair, giving at least \(2^{(a-1)\lfloor b/2\rfloor}\) partitions. Thus \(|A_e|=\exp(o(m))\). If \(D\) is the stabilizer in \(A_m\) on the \(n\)-set, \(A_m=A_eD\) by transitivity, so \(D\) cannot be transitive on the natural \(m\)-set by the same index bound. Nor can it have just orbits of sizes 1 and \(m-1\): the bound on proper transitive subgroups on \(m-1\) would force \(D=A_{m-1}\), contrary to \(n>m\). It thus stabilizes an \(h\)-set, \(2\le h\le m-2\), making \(A_e\) \(h\)-homogeneous on \(m\). This implies 2-homogeneity: the count inside an \(h\)-set of edges from any one pair orbit is constant; comparing sets consisting of \(i\) or \(j\) and \(h-1\) other points shows the adjacency rows at \(i,j\) agree away from both (all such sums of differences vanish). Thus there is just one edge orbit. Even order then gives 2-transitivity by an involution flipping an ordered pair. By the classification of finite 2-transitive groups, alternating simple groups of sufficiently large degree have no such nonnatural action. This proves the bound. The permutation-group inputs used here are the primitive subgroup bound and the classification of finite doubly transitive groups; see \citet{DixonMortimer1996}.
\end{proof}

\begin{lemma}[Changes of defining characteristic]\label{w:characteristic-switch}
For a classical ancestor $i$ and classical descendant $j$ in a prepared path with distinct defining characteristics, there is an absolute $c>0$ such that
\[
 v_i\ge P_j^{c v_j}.
\]
There cannot be three high classical nodes $i<j<l$ with $p_i\ne p_j$ and $p_j\ne p_l$. Consequently, after splitting at alternating nodes and changes of defining characteristic, the path is the union of an absolutely bounded number of consecutive classical runs of one characteristic.
\end{lemma}
\begin{proof}
For high classical pairs in distinct characteristics, absolute constants give \(v_i\ge P_j^{c v_j}\) for ancestor \(i\), descendant \(j\), with \(c>0\). Here is a direct bound. The high classical \(S_j\) contains images of paired special linear groups \({\rm SL}_h(P_j)\), \(h\ge c' v_j\), using a subspace in linear type or paired maximally isotropic coordinates in the form types (one may restrict scalars of matrix entries in the unitary coordinates to \(\mathbf F_{P_j}\)). These act nontrivially including on column shear elements. The restricted projective representation lifts linearly on that \({\rm SL}_h(P_j)\), centrally closed for such high \(h\). Its abelian column shears of size \(P_j^{h-1}\) are diagonalizable here and have some nontrivial character occurring; the \({\rm SL}_{h-1}\) action is transitive on all their nontrivial characters. This proves the assertion. Thus we cannot have classical nodes \(i<j<l\) (indices increasing downwards) with \(p_i\ne p_j,\ p_j\ne p_l\), since then
\[
 \log|S_l|=O(v_l^2\log P_l)=O(\log^2 v_j)
\]
whereas \(v_i\) is at least exponential in \(v_j\). But \(|S_l|\ge v_i^2\) by its absolutely irreducible projective action in \(i\): lift that action to its perfect central cover. Absolute irreducibility implies that its matrices span the full $v_i\times v_i$ matrix algebra. Central lifts of one projective element are scalar multiples, so at most $|S_l|$ matrices are needed in this span. This contradicts the two exponential lower bounds for sufficiently high rank.

Hence there are boundedly many interruptions of consecutive high-classical runs in the \emph{same} characteristic, on paths avoiding bounded types. It suffices to bound those runs. We now work in such a run, notation \(S_i=B_{W_i}\).
\end{proof}

\subsection{Canonical tensor factors over the algebraic closure}

\begin{lemma}[Canonical tensor atoms]\label{w:atoms}
Let $S_i$ be a high classical ancestor, and let $B$ be a candidate simple subgroup of high classical type in the same defining characteristic. Assume that $B$ acts projectively absolutely irreducibly on the ancestor natural module $V$ over $k=\overline{\mathbf F}_p$. Among decompositions of $\operatorname{End}_k(V)$ into nontrivial full matrix factors fixed individually by $B$, there is a unique finest one as a collection of embedded matrix subalgebras. Call its factors the \emph{atoms}. There is at least one atom for each nonzero ordinary $p$-digit in the restricted highest weight of the lifted representation of $\widehat B$. The only possible further separation within a digit is across the multiple bond in types B and C. Every semilinear algebra action, or algebra anti-action with inversion on the group, normalizing the projective $B$-action permutes these atoms.
\end{lemma}
\begin{proof}
To verify this, in any tensor factorization into nontrivial slots with \(B\) fixing the factor algebras individually, lift each projective action to \(\widehat B\). They tensor to the original lift by perfection, and are all irreducible. Regard them as restrictions of rational simples with restricted highest weights for the finite-group parameter of \(B\). Their tensor is then irreducible rationally as well. By Steinberg decomposition, even the subproducts of the contributing digit modules at any one exponent (untwisted) must be irreducible.

Recall the top Levi face of a rational simple \(L(\lambda)\) for nodes \(I\), using weights differing downwards from highest only by \(I\)-roots, is a simple highest module for the Levi derived group with the corresponding coefficients. Indeed it is generated by the highest vector there by distribution-algebra triangular decomposition (PBW); a nonzero Levi submodule on that face has a highest unipotent-fixed vector, fixed also by outside raising groups since they go off the weights. In the simple module this forces the highest line. The face of a tensor product with the sum as highest weight tensors the corresponding faces. Now at any one digit, no two modules from different slots can have active coefficients in the same simple-bond connected component of the Dynkin diagram. Otherwise take a shortest simple-bond path between active nodes of two different slots, possibly just one node (minimize over such pairs in a single simple-bond component). For the two modules the faces on that path have endpoint highest weights \(a\omega_1,b\omega_\ell\) of type \(A_\ell\), \(1\le a,b<p\), with no other supports there by minimality. These are symmetric powers of natural and dual modules (the same natural type for \(A_1\)), simple for these exponents. For example simplicity follows using torus monomials of distinct weights and elementary root transfers of powers with nonzero coefficients at degree less than \(p\). Their tensor is reducible by contraction (one derivative on each side paired), a nonzero equivariant map to a strictly smaller space. This contradicts the face property.

Thus coefficients are disjoint with no carry, and the digitwise sum highest weight is restricted and must be exactly the given tuple for \(V\). Take one atom per nonzero digit except for the following possible refinement in B,C: if the regular chain and the node beyond the multiple bond both contribute and can be separated in a factorization, split them. Indeed in such a separation their two rational pieces tensor irreducibly to the digit module, so all such splits at all digits can be made together. This gives a factorization, and any other one as above has its slot representations given by grouping the atoms (by the coefficient argument and uniqueness of the simples). This is true of the actual matrix subalgebras: equivalent factorizations with fixed slot representation classes have their tensor identifications differing, after intertwiners, by a scalar since \(V\) is absolutely irreducible for \(\widehat B\). In more detail, separated active components at a digit give actual rational factor modules by taking the digit pieces of the slot modules of that separation; their subproduct, untwisted, is irreducible as above, with the individual component weights and their sum all forced. Thus it realizes the indicated split even as rational modules of that digit, independently of the other digits. This proves the claim. Any semilinear or anti-action on the algebra normalizing the projective \(B\)-action (using inverse with an anti-action) permutes the atoms, by the same intrinsic description.
\end{proof}

\begin{lemma}[Descent and the single Frobenius cycle]\label{w:atom-cycle}
Retain the hypotheses of \Cref{w:atoms}, and assume the candidate label has the ancestor tensor-system exclusion in \Cref{w:prepared-path}. Let $\sigma$ be arithmetic Frobenius for the ancestor ground field, except in unitary type where it is the extended semilinear adjoint anti-action with coefficient exponent $P_i$. Then $\sigma$ permutes the atoms in one cycle. If its length is $t=t(B)$, all atoms have a common degree $u\ge2$ and
\[
 u^t=v_i.
\]
The lifts of the candidate normalizer used to act on the endomorphism algebra commute with $\sigma$. In bilateral form types the factors are individually selfadjoint, with alternating factor forms in characteristic two.
\end{lemma}
\begin{proof}
The map $\sigma$ permutes the atoms because it fixes the projective
$B$-action, using inversion with the anti-action in unitary type. We
first specify lifts of the ancestor automorphisms. Use finite-field
projective semilinear representatives, similarities in the form types,
and also graph anti-actions in linear type. Lift their field exponents
by integral powers of split Frobenius on $k$. These lifts commute with
$\sigma$: in unitary type this is the hermitian-similarity identity
extended to $k$, and otherwise it follows from the finite-field matrices.

The finite projective representatives are faithful on the ancestor
simple group. Indeed, if a collineation with coefficient exponent
$0\le a<\log_p|F|$ centralizes it, perfection identifies its natural
finite lift with the coefficient twist of highest weight $p^a\omega_1$.
Restricted-weight uniqueness forces $a\equiv0\pmod{f_i}$. In unitary
type it also rules out $a=f_i$, since wrapping produces the distinct
diagram-dual natural weight. With the field part removed, absolute
irreducibility leaves only a scalar. The same weight comparison rules
out a linear-type duality in the kernel. Thus closure lifts differ only
by arithmetic ground-field Frobenius powers. A system stable under those
powers as a collection has a well-defined full stabilizer in the finite
ancestor. In unitary type the representatives here are collineations;
$\sigma$ itself is the semilinear adjoint anti-action and $\sigma^2$ is
arithmetic Frobenius for $F$.

We next check the forms needed for descent. In bilateral types the
perfect lift of $B$ preserves the ambient form. Restricted-weight
uniqueness makes its digits self-dual, and any pieces separated across
the B or C multiple bond are self-dual as well. Their invariant forms
tensor to the ambient form by uniqueness, so the atom algebras are
individually selfadjoint. In characteristic two each factor form is
alternating: a nonalternating symmetric invariant form would give the
nonzero invariant functional $x\mapsto\sqrt{b(x,x)}$ on a nontrivial
irreducible module.

If a quadratic form is present in characteristic two, every grouping
into at least two aggregate factors has the canonical quadratic with
the given polar form and value zero on simple tensors. To construct it,
prescribe zero on tensor-basis vectors. On a simple tensor, polar terms
with a repeated index vanish and the remaining terms cancel in pairs
by exchanging one index. This quadratic is invariant under the lifting
group. Its difference from the given invariant quadratic is the square
of an invariant linear functional, hence zero.

Suppose $\sigma$ has more than one orbit on atoms, and tensor together
the atoms in each orbit. The aggregate algebras descend under arithmetic
Frobenius to central simple $F$-factors, which split over the finite
field. In unitary type first descend under $\sigma^2$; stability under
$\sigma$ then gives selfadjointness for the hermitian adjoint. Bilateral
selfadjointness was just proved. Candidate lifts commute with $\sigma$
and therefore permute these aggregate factors. This is exactly a
ground-field tensor system forbidden by the preceding exclusions.

Its full stabilizer is proper. Otherwise $S_i$ acts trivially on at most
$\log_2v_i$ slots, by simplicity and comparison of its order with the
corresponding factorial. Absolute irreducibility makes each slot action
nontrivial. A smallest slot has degree at most $\sqrt{v_i}$, contrary to
\Cref{w:min-degree} at the fixed high threshold. There is therefore one
$\sigma$-cycle. Its atoms have equal degree $u$, and $u^t=v_i$.
\end{proof}

\subsection{Counting proper containing stabilizers}

The next two counts supply the finite lists needed for selection. Each
full stabilizer contains the candidate label, so it is an actual label
in the ancestor interval. To count distinct primes, we construct a
subgroup lying in all previously imposed stabilizers but not in the
next one. Their intersections then descend strictly, and interval
height bounds the number of such steps. Tensor coarsenings leave a
bounded small-slot remainder; fixing roots do not.

\begin{proposition}[Counting prime tensor coarsenings]\label{w:coarsening-count}
Under the hypotheses of \Cref{w:atom-cycle}, each prime $s\mid t(B)$ defines a canonical congruence coarsening $\mathcal P_s$ into $s$ matrix factors, normalized by the candidate label. Its full stabilizer in the ancestor is proper. There is an absolute constant $c_1$ such that
\[
 \#\{s:s\text{ is prime and }s\mid t(B)\}
 \le \operatorname{ht}[K_{\rm cand},K_{W_i}]+c_1.
\]
The height is taken in the actual subgroup interval, or equivalently in its prepared ancestor quotient.
\end{proposition}
\begin{proof}
For each prime $s\mid t$, group the atoms by their residue classes modulo
$s$ around the $\sigma$-cycle. Candidate lifts commute with the cycle,
so they normalize this canonical coarsening $\mathcal P_s$. Its full
stabilizer is well defined by \Cref{w:atom-cycle}. If the entire socle
$S_i$ preserved it, its permutation action on the $s$ slots would
commute with an $s$-cycle. That image is cyclic, hence trivial by
perfection. A slot of degree $u^{t/s}\le\sqrt{v_i}$ would then carry a
nontrivial projective action, contradicting \Cref{w:min-degree}.
Thus each coarsening gives a proper containing label.

Let $d$ be the product of distinct primes already imposed. Their common
refinement is $\mathcal P_d$, with $d$ slots of degree $h=u^{t/d}$.
In $G_i(k)$ take the product $X_d$ of the supported connected projective
groups on these slots: $\operatorname{PGL}_h$ in ambient type A, and
the connected projective classical groups of the factor forms in
bilateral types. Put $X_1=G_i(k)$ itself. In particular, at $d=1$ we
keep the orthogonal ambient group in characteristic two. For $d>1$
the tensor construction embeds these products in the connected ambient
group; the canonical-quadratic observation in \Cref{w:atom-cycle}
justifies the characteristic-two orthogonal case.

The map $\Phi_i$ cycles the $d$ factors through $\sigma$, using inversion
on units in the anti-case. On one factor $\Phi_i^d$ is an ordinary
positive Steinberg map of parameter $P_i^d$: after this coefficient
twist, its matrix-algebra action is an isomorphism or anti-isomorphism
preserving the factor group. The fixed group $X_d^{\Phi_i}$ therefore
projects isomorphically to that one-factor fixed group. While $h$ is
sufficiently large, its derived group gives a full high-rank natural
simple $Y$ on the chosen slot. It lies in $S_i$ by taking commutators
and preserves every previously imposed coarsening.

For a new prime $s$, suppose $Y$ also preserves $\mathcal P_s$.
Perfection and commutation with the induced $s$-cycle make it fix those
slots individually. Inside the chosen $d$-slot, the two coarsenings
intersect in $s$ tensor factors of degree $h^{1/s}$, by the Chinese
remainder description of their atom indices. At least one factor
action is projectively nontrivial, since their tensor is the natural
action of $Y$ on that slot. Its degree contradicts \Cref{w:min-degree}.
Thus the next stabilizer makes the intersection strictly smaller.

Interval height bounds the primes added while $h$ is high. At stopping,
$h\ge2^{t/d}$ and $t/d$ is divisible by every unused distinct prime.
Bounded $h$ leaves only absolutely many unused primes. This gives the
claimed additive constant $c_1$ in the actual interval, whose height is
unchanged by the centralizer quotient.
\end{proof}

\begin{proposition}[Counting fixing prime roots]\label{w:root-count}
Let $B$ be as in \Cref{w:atom-cycle}, in the high classical ancestor with adjoint group $G_i$ and ordinary Steinberg map $\Phi_i$ of height $f_i$. Consider pairs $(s,\Theta)$ where $s$ is prime,
\[
 \Theta^s=\Phi_i,
 \qquad \Theta|_B=1,
\]
and $\Theta$ belongs to the inner--diagram--Frobenius point-map group. The full normalizer labels of $S(\Theta)=(G_i^\Theta)'$ are proper, contain the candidate label, and number at most
\[
 2\operatorname{ht}[K_{\rm cand},K_{W_i}].
\]
There are at most two maps at each possible height, one for each diagram choice, and all such roots commute.
\end{proposition}
\begin{proof}
The inner centralizer of $B$ is trivial by projective absolute
irreducibility and perfection. Modulo inner automorphisms the point-map
group is abelian, with at most two diagram choices at any height.
Consequently two maps of the same height and diagram choice fixing $B$
agree, and all fixing roots commute. Conjugating a root by a candidate
lift changes no outer class and still fixes $B$ pointwise. The same
argument shows that the lift commutes with the root and normalizes
$S(\Theta)$.

An inner element normalizing $S(\Theta)$ likewise commutes with $\Theta$,
by applying this centralizer argument to the natural absolutely
irreducible fixed simple. Since $G_i$ is adjoint, that element itself
lies in $G_i^\Theta$. The root has height $f_i/s<f_i$, so the order
comparison makes its normalizer proper in the ancestor. We have thus
obtained actual proper labels containing the candidate.

Choose one root $\Theta_s$ for each occurring prime. For a product $d$
of primes already chosen, select integers $a_s$ with
$\sum_{s\mid d}a_s d/s=1$ and put
\[
 \Theta_d=\prod_{s\mid d}\Theta_s^{a_s}.
\]
The roots commute, and $\Theta_s^s=\Phi_i$ gives
\[
 \Theta_d^{d/s}=\Theta_s,\qquad
 \operatorname{ht}(\Theta_d)=f_i/d.
\]
For example, raising the product to $d/s_0$ gives
$\Theta_{s_0}^{a_{s_0}d/s_0}$ times
$\Phi_i^{\sum_{s\ne s_0}a_s d/(ss_0)}$, which is exactly
$\Theta_{s_0}$. This verifies the identity also with diagram components.
For the empty product set $\Theta_1=\Phi_i$.

The group $S(\Theta_d)$ is fixed by each previous root and lies in all
their normalizer labels. If it normalized the fixed simple for a new
prime root $\Theta_l$, the preceding inner-normalizer argument would
make it fixed by $\Theta_l$ as well. It would then lie in the fixed group
of the combined root $\Theta_{dl}$, contradicting
\[
 |S(\Theta_d)|>|G_i^{\Theta_{dl}}|.
\]
Each new prime therefore makes the label intersection strictly smaller.
There are at most as many primes as the interval height, and at most
two roots at each associated height. The number of normalizer labels
is consequently at most twice that height.
\end{proof}

\begin{corollary}[Equal natural modules force a fixing root]\label{w:natural-root}
Suppose $B<S_i$ is a same-series high classical candidate of the same natural dimension as $S_i$, and its projective action on $V_i$ is its own natural module, allowing a Frobenius twist and, in type A, a diagram natural/dual twist. Then some prime root of $\Phi_i$ fixes $B$ pointwise. Thus the absence of such fixing roots excludes this case.
\end{corollary}
\begin{proof}
We record why root-free protection precludes equal-natural embeddings within one root series. If the action of \(B\) on \(V_i\) is projectively its own natural module (allowing Frobenius twists, also the diagram natural/dual twist in A), both groups in the same classical series and dimension, then \(B=S(\Psi)\) inside \(G_i\) for a standard Steinberg map up to conjugacy. Indeed identify the natural spaces over \(k\) by the representation equivalence (split Frobenius twists preserve the standard finite images); forms match up to scalar by uniqueness, also the quadratic when needed by square-difference uniqueness on the finite group. This identifies the corresponding algebraic adjoint groups. Divide the height of \(\Phi_i\) by that of \(\Psi\); after composing off that power of \(\Psi\) we still fix \(B\) pointwise, thus cannot have positive remainder by the order comparison. At height zero, an inner centralization is trivial; so is a possible D-diagram centralization (acting in \({\rm PGL}(V_i)\)); and a non-inner A-graph centralization would give projective self-duality of the finite natural representation, impossible by the distinct restricted weights and perfection. So \(\Phi_i=\Psi^z\), \(z>1\) by strict inclusion, and some prime root fixes \(B\).
\end{proof}

\subsection{Simultaneous protection in a pruned chart}

\begin{definition}[A protected classical run]\label{w:protected}
A consecutive high classical run of one defining characteristic is \emph{protected} if for every ancestor--descendant pair within the run the descendant has one tensor atom on the ancestor natural module and no prime root of the ancestor Steinberg map fixes it pointwise. The path also retains all properties of \Cref{w:prepared-path}.
\end{definition}

\begin{proposition}[Selecting all ancestor protections]\label{w:selection}
Assume that the current same-characteristic classical run is already protected through its current node $W$. Fix an upper bound $D$ on the number of path steps and a fixed number $c$ of continuation colors. There is a finite extension rank $r$, depending only on $D,c$ and earlier fixed requirements, such that from every suitably pruned private chart at a current high classical node one can select a next proper extension. If that extension continues the current same-characteristic run, it has both protections relative to every ancestor in that run, including the current node. Earlier structural exclusions are preserved.
\end{proposition}
\begin{proof}
Work inside a successful pruned flat of extension rank $r$ at the current
node $W$. All its comparable candidate labels have the prepared
properties and actual nested simples, by \Cref{sec:patterns}.
Color its proper extensions with at most $c$ colors, recording in
particular whether their socle continues the current same-characteristic
high-classical run. Consider a uniform random flag at extension ranks
\[
 \ell,2\ell,\ldots,(c+1)\ell,\qquad r=(c+1)\ell.
\]
Every flag has an equal-colored pair. Hence some two ranks $a<b$ have
matching probability at least $\binom{c+1}{2}^{-1}$. Conditioning on the
$b$-extension gives a particular $U$ for which at least that fraction
of the $a$-extensions $L\subset U$ have its color. If this color does
not continue the run, any such $L$ needs no additional protection.

Otherwise apply the two counts to $B_U$ in each ancestor $i$ of the
current run, including $W$ itself. These lists detect every failure of
protection at a matching $L$. Indeed, if $t(B_L)>1$, a prime congruence
coarsening of its atoms has a stabilizer containing $K_L$. Since
$B_U\le B_L$ fixes its factors individually, each factor groups atoms
of $B_U$. The partition is invariant under the single Frobenius cycle;
a cycle has a unique invariant partition into $s$ congruence classes.
It is therefore one of the coarsenings counted for $B_U$. A prime root
fixing $B_L$ likewise fixes $B_U$ and occurs on its root list.

None of the counted labels contains $K_W$. For an earlier ancestor,
such containment would make $B_W$ fixed by the root, or individually
fix the coarsening slots: its slot permutation commutes with the cycle
and is trivial by perfection. Either conclusion contradicts the
protection already imposed on $W$. For the current ancestor, properness
of the counted labels gives the same noncontainment. Thus each bad label
is $K_Z$ for a subflat $Z\subseteq U$ with $Z\not\subseteq W$.

For a uniform $a$-extension $L$ in $U$, the probability of $Z\subseteq L$
is at most $2q^{a-b}$: it suffices to require one fixed point of
$Z\setminus W$ in the ordinary affine-flat count. Here $q$ is the chart
field size. The two counts give $O(1+\operatorname{ht}[K_U,K_{W_i}])$
bad labels per ancestor. Reverse rank differences bound this height by
$Dr$ over at most $D$ path steps. The total number is therefore at most
$CD(1+Dr)$ for an absolute $C$ after the colors and thresholds are fixed.
Since $b-a\ge\ell$, their union occupies a fraction at most
\[
 2CD(1+D(c+1)\ell)q^{-\ell}
\]
of those extensions. This tends to zero with $\ell$, whereas the
matching-color fraction is at least $\binom{c+1}{2}^{-1}$. Choose $\ell$
large enough also for the earlier fixed requirements. A matching $L$
outside the union has one atom and no fixing prime root relative to
every ancestor in the run. It retains the structural exclusions because
it lies in the successful pruned flat.

There is no circular choice of parameters: the following path bound
concerns paths already satisfying these group protections. Its absolute
constant is obtained first, then $D$ is fixed, and only then are $\ell$
and $r$ chosen.
\end{proof}

\subsection{From finite inclusions to rational compositions}

The tensor protection makes each edge a single ordinary $p$-digit.
To compose two edges rationally, the lower algebraic group must preserve
the middle group's defining form. Bilinear forms follow from
restricted-weight uniqueness. Quadratic forms in characteristic two
require the next two lemmas. Root protection will have a separate role:
it excludes natural coefficient-one edges by \Cref{w:natural-root}.

\begin{lemma}[Restriction is injective in degree-one cohomology]\label{w:h1-injective}
Let $J$ be a simply connected simple classical algebraic group in the
high-rank range of \Cref{w:classical-conventions}, over
$k=\overline{\mathbf F}_p$. Let
$\Phi=\tau\operatorname{Fr}_p^f$ be an ordinary pinned diagram--field
Steinberg map with $f\ge1$, parameter $P=p^f$, and diagram action $\tau$
on weights. If $D'=L_J(\xi)$ is a nontrivial $P$-restricted simple
rational $J$-module, then
\[
 H^1(J,D')\longrightarrow H^1(J^\Phi,D')
\]
is injective. The finite group here is the full simply connected fixed group.
\end{lemma}
\begin{proof}
Put $\mathcal G=k[J]^{J^\Phi}$, with right finite invariance and left
translation. It suffices to show
\[
 (D'\otimes(\mathcal G/k))^J=0.
\]
Indeed, a finite-fixed vector splitting an extension of $k$ by $D'$
defines the equivariant function sending $g$ to its translate of that
vector. It is right finite invariant, and its image in the quotient
module is constant. Modulo constant functions it therefore lies in the
displayed invariant space. Vanishing makes the function itself constant,
giving a rationally fixed lift.

The Lang map $g\mapsto g\Phi(g)^{-1}$ identifies the right quotient
$J/J^\Phi$ with $J$ \citep{Lang1956,Steinberg1968}
\citep[Definition~17.95 and Theorem~17.96]{Milne2022}.
It is surjective and étale since $d\Phi=0$, hence faithfully flat.
Its fibres are exactly the right cosets of the finite reduced
fixed-point scheme: $h=gx$ with $x=\Phi(x)$. Descent of regular functions
identifies $\mathcal G$ with the regular function algebra carrying the
twisted left-right action.

The Donkin--Koppinen filtration begins with constants and has tensor
pairs of induced highest-weight modules as its sections
\citep[Theorem~4.3 and its preceding construction]{Mathieu2000}
\citep[Section~2.4]{BendelNakanoPillen2011}. Pulling this two-sided
filtration back to the twisted diagonal gives an exhaustive filtration
of $\mathcal G/k$ with sections
\[
 \nabla(\nu)\otimes\nabla(\nu^*)^{[\Phi]},
 \qquad \nu\ne0\ \text{dominant}.
\]
Here $*$ denotes highest-weight dual, $[\Phi]$ pullback, and $\Delta$
below is a Weyl module. A nonzero invariant after tensoring a section
by $D'$ gives a nonzero map
\[
 \Delta(\nu)^{[\Phi]}\longrightarrow L(\xi)\otimes\nabla(\nu).
\]
The source is still generated by its highest vector, since $\Phi$ is
surjective on algebraic-group points. The map is nonzero there, so
\[
 P\tau(\nu)\le\xi+\nu
\]
in dominance order. The source is also trivial on the $f$-th Frobenius
kernel $J_f$. Its nonzero image therefore gives
$\operatorname{Hom}_{J_f}(L(\xi)^*,\nabla(\nu))\ne0$, with kernel
and invariance understood as group schemes.

Since $\xi$ is $P$-restricted, the restricted-simple theorem and
Steinberg decomposition force a rational composition factor of
$\nabla(\nu)$ whose highest weight has the form
\[
 \xi^*+P\gamma\le\nu,\qquad \gamma\text{ dominant}.
\]
Indeed, every rational simple restricts to copies of the $J_f$-simple
for its $P$-restricted part, and those restricted simples are pairwise
distinct. Pairing the two inequalities with $\rho^\vee$ gives the
diagram- and dual-invariant height comparisons
\[
 (P-1)\operatorname{ht}(\nu)
 \le\operatorname{ht}(\xi)\le\operatorname{ht}(\nu).
\]

Only $P=2$ and equality remain. Then $p=2$, $f=1$, and $\xi$ has an
active coefficient one. Equality gives $P\tau(\nu)=\xi+\nu$, so the
image of the source highest vector lies on the unique tensor highest
line. Let $v_\xi\otimes v_\nu$ span it and choose $\alpha$ with
$\langle\xi,\alpha^\vee\rangle=1$. The lowering operator has value
\[
 (f_\alpha v_\xi)\otimes v_\nu
   +v_\xi\otimes(f_\alpha v_\nu).
\]
The first summand is nonzero by the simple-root bracket, and the
first-factor weights $\xi-\alpha$ and $\xi$ differ. The summands cannot
cancel, including in characteristic two. The line is not $J_1$-fixed,
contradicting the kernel-trivial source action. Thus no section has
an invariant. Every invariant would occur in the exhaustive filtration
and have a nonzero successive image, so the required vanishing follows.

The restricted-kernel and regular-bimodule results are the simply
connected theorems in \citet[Part II, Chapter 3 and II.4.20]{Jantzen2003};
see also \citet[Section~1.4]{LinNakano2006} and
\citet[Theorem~4.3]{Mathieu2000}. The pullback to the ordinary
diagram-twisted diagonal uses no complete reducibility.
\end{proof}

\begin{lemma}[The quadratic form in characteristic two]\label{w:quadratic}
Let $J$ be simply connected of sufficiently high type $D_e$ in characteristic two, with ordinary pinned map $\Phi=\tau\operatorname{Fr}_2^f$, $f\ge1$. Let $\eta\ne0$ be $2$-restricted, and consider $L_J(2^y\eta)$. Suppose a quadratic form $Q$ on this module is invariant under the full finite cover $J^\Phi$, and its polar form is rationally $J$-invariant. Then $Q$ itself is rationally $J$-invariant.
\end{lemma}
\begin{proof}
Untwist $2^y$ to $L_J(\eta)$; split Frobenius permutes $J^\Phi$
because $\Phi$ is a standard pinned graph/field map. Exact invariance
under the chosen finite linear lift holds by perfection: its similitude
scalar is a character. Rational invariance of the polar form makes
\[
 c(g)=gQ-Q
\]
a regular cocycle with values in the subspace of squares of linear
forms. We will take its coefficientwise square root to obtain a cocycle
in $L_J(\eta)^*$.

The target Lie action on squares is zero, so the differential of $c$
at the identity kills commutators. The Lie algebra of simply connected
type D is perfect also in characteristic two: simple coroots span the
torus part and are opposite-root brackets, and every D-root pairs oddly
with some coroot. Thus this differential, and by translation every
differential of $c$, vanishes.

Each coordinate of $c$ is consequently a square in the function field,
using the kernel of differentials over a perfect field. Its square root
is regular by normality. Taking roots of the coefficients in a fixed
dual basis therefore gives a regular linear-form cocycle, zero on
$J^\Phi$. Additivity, injectivity and equivariance of squaring on
$k$-points give its cocycle identity.

The square-root cocycle takes values in the nontrivial restricted dual module after untwisting. By \Cref{w:h1-injective}, its class is zero, so it has the form $g\ell-\ell$ for a linear functional $\ell$. Its vanishing on $J^\Phi$ makes $\ell$ finite-fixed. The restricted finite module is nontrivial and irreducible, so $\ell=0$. Thus the square-root cocycle, and therefore $gQ-Q$, is identically zero. The perfect-Lie-algebra argument above is used here only for type D in characteristic two.
\end{proof}

\begin{proposition}[Composing the rational highest-weight actions]\label{w:rational-composition}
Let $i<j<l$ lie in one protected run, and suppose these three nodes have the same series A, B, C, or D. Use consistent natural realizations and simply connected covers of the actual nested simple subgroups. Write $G_s=\widehat G_j$ and $J=\widehat G_l$. The lifted edge representations $ij$ and $jl$ have highest weights
\[
 p^x\lambda,\qquad p^y\eta,
\]
where $\lambda,\eta\ne0$ are $p$-restricted. The latter representation induces an algebraic homomorphism $J\to (G_s)_{\rm ad}$ which lifts to $J\to G_s$. The resulting rational composition of $L_{G_s}(p^x\lambda)$ is irreducible. Its highest weight $\mu$ is one nonzero $p$-digit times a power of $p$, and that digit has the shape, up to diagram, of the direct edge $il$.
\end{proposition}
\begin{proof}
Consider path nodes \(i<j<l\) in one protected run, of the same root series A, B, C or D. Use consistent finite-group identifications at each node (actual simple \(B_{W_j}\), its projective natural realization, and simply connected cover, for both its incoming and outgoing edges). Thus the projective restriction along \(i\to j\to l\) is just via inclusion of the same nested simples as \(i\to l\); the weight colors below also respect diagram changes. Write \(G_s=\widehat G_j\) and \(J=\widehat G_l\) temporarily for the simply connected algebraic groups. Choose pinned coordinates with standard graph/field Frobenius on each (conjugating the finite realizations accordingly if necessary). The edges \(ij,jl\), evaluating natural modules on descendants projectively, have highest weights \(p^x\lambda,p^y\eta\) respectively for the lifts, \(\lambda,\eta\ne0\) \(p\)-restricted, since each has one atom. We claim \(L_J(p^y\eta)\) on the natural space for \(j\) gives a morphism \(J\to (G_s)_{\rm ad}\) inducing the given projective action on finite points. This is immediate in A. In the remaining types the finite representation preserves the bilateral form (perfection); by restricted-weight uniqueness it is rationally self-dual, and the invariant form rationally therefore agrees up to scalars. This suffices except for D in characteristic two, where \Cref{w:quadratic} supplies rational invariance of the quadratic form as well.

The morphism to the adjoint group (in form types one can first map by the given linear representation to the connected classical isometry group, using connectedness of \(J\), then take its adjoint quotient) lifts to \(G_s\) by simply-connectedness. For example pull back the central isogeny; the reduced identity group in the pullback maps surjectively to \(J\) (finite surjection, finitely many components), by a central isogeny with kernel of multiplicative type; its radical has trivial image and hence dimension zero, so it is semisimple and maps isomorphically by the simply connected root datum of \(J\). Its natural action agrees linearly with the above one (no scalar characters). On finite points the lift differs at most by central elements from lifts to the usual finite covering group at node \(j\). Indeed both give the prescribed projective matrices on the natural space (the image of the finite cover surjects projectively onto the given socle), and the kernel of the projective natural map on simply connected algebraic-group points is central. Thus composing \(L_{G_s}(p^x\lambda)\) yields on restriction to \(J^\Phi\) projectively the direct-edge action \(il\). The rational composition is irreducible by finite irreducibility. Its highest weight \(\mu\) for \(J\) is again a single \(p\)-digit times a power of \(p\), by Steinberg decomposition and the protected individual tensor indecomposability on finite points. Indeed multiple rational digits here would give multiple nontrivial-dimensional factors fixed individually by the finite cover, hence also as matrix factor algebras by the finite projective \(S_l\)-action (and this projective action already agrees as just explained). Moreover the digit shape agrees up to diagram with the direct edge's: reduce its Frobenius twist exponent on the finite group using the defining graph/field map and apply restricted-weight uniqueness (linear lifts agree by perfection).
\end{proof}

\subsection{Universal cuts of the full weight multiset}

\begin{lemma}[Chamber maxima and universal cuts]\label{w:universal-cuts}
Use the triple, groups, and weights of \Cref{w:rational-composition}. Let $\Omega$ be the full weight multiset of $L_J(\eta)$, with every multiplicity retained. After mapping a maximal torus of $J$ into one of $G_s$, the natural weights restrict to $p^y\Omega$. One Weyl translate of $\lambda$ maximizes throughout the positive cocharacter chamber of $J$, and on each active fundamental summand of $\lambda$.

In type A an active node $a$ of the upper weight $\lambda$ therefore
gives a cut selecting $a$ entries of $\Omega$, each dominating every
unselected entry, with no equal-vector tie across the cut. In B and C
the analogous signed fundamental cut selects $a$ entries, including
the last node up to the B spin scale; its selected weights are nonzero
and in the positive cone. This also holds at the nonspin D nodes when
$\lambda$ has no spin support. At an active D spin node the maximum
instead chooses one orientation $u_a$ of each natural signed pair, and
\[
 u_a+u_b\text{ is nonnegative on the chamber whenever }a\ne b.
\]
Distinct ambient orbit weights cannot restrict to the same highest weight. More generally, every weight of a simple highest module differing from highest by a multiple of one simple root has multiplicity at most one.
\end{lemma}
\begin{proof}
Map a maximal torus of $J$ into one of $G_s$. The restricted natural
multiset is $p^y\Omega$. By irreducibility of the composition and the
highest-weight convex hull property, the maximum pairing on the orbit
of $p^x\lambda$ is the linear function given by $\mu$ throughout the
positive cocharacter chamber of $J$. Choose a maximizing orbit weight
at a strictly interior cocharacter. Its restriction must equal $\mu$
as a vector, so that same Weyl translate $w\lambda$ maximizes everywhere
in the chamber.

It also maximizes each active fundamental summand. For any one
cocharacter a common Weyl translate maximizes all fundamental summands;
equality in their positive sum therefore forces equality for each.
Two distinct ambient orbit weights cannot both restrict to the
composed highest weight, whose multiplicity is one.

Suppress the positive scales $p^x,p^y$ for the cut description. In type
A an active upper node $a$ selects $a$ natural coordinates at $w\lambda$.
Exchanging one selected and one unselected coordinate shows, throughout
the lower chamber, that every selected restriction dominates every
unselected restriction. An equal-vector tie would give another
maximizing ambient orbit weight across a strict coefficient boundary,
contradicting highest multiplicity one.

In B and C the analogous choice takes $a$ entries of the full signed
multiset, also at the last node up to the B spin scale. In D the same
argument applies before the spin nodes when $\lambda$ has no spin
support. Sign changes show that selected restrictions lie in the
positive cone; in D compensate by an unused coordinate. A zero-vector
restriction would permit a sign change in the full maximizing orbit
weight, using a zero-coefficient coordinate for D parity, and again
contradict multiplicity one. Thus the selected values are nonzero.
Signed coordinate exchanges compare them with every unused value,
including negative choices and the extra zero in B. An equal nonzero
tie across the cut also changes the ambient orbit weight; parity can
be preserved by changing both directed signs when needed. These cuts
are therefore strict as asserted.

For D spin support, take an active spin fundamental summand instead.
Its maximizing signed choices $u_i$ satisfy the stated pair-sum
condition, because flipping any two choices cannot increase the
pairing on the chamber.

These comparisons concern all entries of $\Omega$, including nonextremal weights and repeated entries. A cut of its Weyl orbit alone is only a necessary consequence. The single-root multiplicity assertion follows because the corresponding weight space of the highest-weight module is generated by the one divided-power monomial in that root; it does not assert multiplicity one for arbitrary weights.
\end{proof}

\subsection{The word argument in type A}

\begin{lemma}[Type A orbit cuts]\label{w:a-word-cuts}
Let $J$ have sufficiently high type $A_{e-1}$, and let $0\ne\eta$ be $p$-restricted. Suppose a nonempty proper universal cut of the full weight multiset of $L_J(\eta)$ is given as in \Cref{w:universal-cuts}. Replace the cut by the longest-element image of its complement if necessary so its selected part of the extremal orbit has size at most half that orbit. Unless $\eta$ is supported at one endpoint, the selected orbit words consist of the highest word alone, possibly its single boundary swap for a binary word, or the positive half of the orbit of the exact highest root. In the final case the highest weight is exactly $\epsilon_1-\epsilon_e$, not a larger restricted multiple.
\end{lemma}
\begin{proof}
Take \(J\) of type \(A_{e-1}\). Its orbit on \(\eta\) consists of permutations of a descending nonconstant word, ordered for dominance by prefix sums (coordinate sum zero). In a universal cut of \(\Omega\), the upper orbit words form a nonempty proper subset \(U\), since the orbit contains the full highest and lowest weights. We may study the smaller orbit side, replacing the full cut by the longest-element reversal of its complement when needed (this is again an upper cut of the weight multiset, now of complementary size) so \(|U|\) on the orbit is at most half. Then \(U\) is an initial segment both for lexicographic decreasing order and for lexicographic increasing order starting from the right. This follows by the first differing coordinate comparisons from strict cone domination and equality of total sums.

Here is the description needed of such segments. Let \(M\) be the largest letter, multiplicity \(b\), and \(m\) the smallest, multiplicity \(c\). If \(U\) contains a word not starting with \(M\), it contains all starting with \(M\). If \(b\ge2\) this includes one also ending with \(M\), forcing all ending with lower letters; the two fractions (and extra words) contradict the half bound. Thus \(b=1\). With second-largest letter \(T\), a starting-\(M\) word ending in \(T\) forces in all ending below \(T\). If there are at least two such lower letters counting multiplicity, these include one also starting below \(T\), forcing all starting at or above \(T\). Again the fractions contradict the bound. Namely if \(l\) letters are below \(T\), the proportions ending below \(T\) and starting at or above \(T\) are \(l/e\) and \(1-l/e\); in the indicated situation \(U\) contains each such class entirely and strictly. We obtain only a natural two-value shape with singleton extreme against constant, or \(M,T^{e-2},m\). Symmetrically these are the exceptions if some selected word does not end with \(m\).

Otherwise all words in $U$ start with $M$ and end with $m$. Until the
first word not ending with $m$, left lex order fixes the first $e-c-1$
coordinates of the descending word: the greatest word not so ending
already has those coordinates. Right lex order confines changes to
the first $b+1$. Only positions in $[e-c,b+1]$ can therefore change.
This interval has at most one position unless $b+c=e$, the binary
case; then its two positions permit only the single boundary swap.
Thus, outside the exceptional shapes, $U$ consists of $\eta$ alone or
of $\eta$ and that swap.

For \(M,T^{e-2},m\), use positions \(a,d\) of the largest and smallest letters. Both orders begin with the half \(a<d\), one prioritizing \(a\) increasing there and the other \(d\) decreasing. A proper nonempty common prefix of the half has only \((1,e)\) (taking \(a>1\) forces \((1,2)\), then the full half; similarly for \(d<e\)). Dominance across the cut consisting of the half requires \(M-T=T-m\), by comparing words with both defects early and both defects late, in both choices of early side. At a prefix separating the two defects of one word (both before the boundary) from those of the other (both after) this compares \((M-T)+(m-T)\) to zero, in both directions. Thus \(\eta=z(\epsilon_1-\epsilon_e)\) in that case (\(\epsilon_i\) standard coordinates). In the full weight set each selected weight then has all prefix sums nonnegative (dominance over negative scaled simple roots on the other side), and each unselected one has them nonpositive. If \(z>1\), the weight \(\eta-\alpha_1\) occurs by simple-root lowering (\(z<p,\ \alpha_1=\epsilon_1-\epsilon_2\)), and can be permuted to mixed prefix signs. So this half exception requires \(z=1\).
\end{proof}

For edges whose descendant has type A, color the nonzero restricted digit by the following five possibilities:
\begin{enumerate}
\item one active endpoint with coefficient one;
\item one active endpoint with coefficient greater than one;
\item one active interior node;
\item multiple active nodes, excluding the exact highest root;
\item the exact highest root.
\end{enumerate}
This palette is invariant under the A diagram reversal.

\begin{proposition}[No homogeneous type A configuration]\label{w:a-obstruction}
Four nodes of a protected run of type A cannot have all pair edges of one color in the preceding palette.
\end{proposition}
\begin{proof}
Take the first three of four homogeneous nodes, with the notation of
\Cref{w:rational-composition}. For each color we first constrain the
extremal orbit cut, then use actual full weights or their multiplicities.

\paragraph{Multiple support and the highest root.}
In the multiple nonexact color, the orbit lemma leaves only $\eta$ on
the smaller side of every cut. A full selected weight $\nu$ must
dominate every active simple reflection $s_h\eta$. Thus the nonnegative
difference $\eta-\nu$ is supported only at $h$ for each active $h$.
Multiple support forces $\nu=\eta$, whose multiplicity is one. Every
active node of the upper digit $\lambda$ must consequently be an
endpoint; since it has multiple support, both endpoints are active.
Apply the same argument to the last three nodes to obtain that
conclusion for $\eta$ as well. In the exact highest-root color both
digits already have these two active endpoints.

Choose natural coordinates at a maximum of $\lambda$ whose first,
second, penultimate and last restrictions, before the lower twist, are
\[
 \eta,\quad \eta-\alpha_1,\quad
 w_0\eta+\alpha_1,\quad w_0\eta,
\]
respectively. Endpoint activity supplies these distinct weights in
the lower module. The middle coefficients of $\lambda$ are constant,
so their placement preserves the maximum. Write $\theta=w\lambda$
for this maximizing weight, and $\beta_1,\beta_{v_j-1}$ for the two
endpoint simple roots in that coordinate order. Restricted activity
makes $\theta-\beta_1$ and $\theta-\beta_{v_j-1}$ distinct weights of
the upper module. Their restrictions satisfy
\[
 \operatorname{res}\bigl(p^x(\theta-\beta_1)\bigr)
 =\operatorname{res}\bigl(p^x(\theta-\beta_{v_j-1})\bigr)
 =\mu-p^{x+y}\alpha_1.
\]
The two ambient weight spaces stay independent on restriction. They
give multiplicity at least two along one simple-root ray from the
composed highest weight, contrary to \Cref{w:universal-cuts}.

\paragraph{One interior node.}
Write $\lambda=c\omega_a$ and $\eta=z\omega_h$, with $a$ and $h$
interior in their respective diagrams. Replace the full cut by the
longest-element image of its complement if required by
\Cref{w:a-word-cuts}. Its selected orbit words are $\eta$, possibly
together with $s_h\eta$.

Every full selected weight $\nu$ differs from $\eta$ only in the
$\alpha_h$ direction. If $s_h\eta$ is outside, domination of that
word proves this. Otherwise use the two outside words
$s_{h-1}s_h\eta$ and $s_{h+1}s_h\eta$: their downward supports intersect
only at $h$. The actual module also contains
\[
 \zeta=\eta-\alpha_h-\alpha_{h-1}.
\]
To obtain it, lower at $h$ and then across the simple bond to its left,
where the new coefficient is one. This weight is outside the cut,
because its downward support is not confined to $h$. Domination of
$\zeta$ leaves only $\eta$ and $\eta-\alpha_h$ on the selected side,
each of multiplicity at most one.

The cut has $a$ entries, or $v_j-a$ in the complemented case. Both
numbers are at least two, so these are exactly its two entries.
Since the full multiset $\Omega$ sums to zero, complementation and
the longest element replace a sum by its highest-weight dual. Hence
\[
 \mu=c\,p^{x+y}(2\eta-\alpha_h)
 \quad\hbox{or}\quad
 \mu=c\,p^{x+y}(2\eta-\alpha_h)^*.
\]
Both neighboring fundamental coefficients of $2\eta-\alpha_h$ are
positive. No common Frobenius rescaling removes either support node,
contradicting the one-node digit shape of the direct edge.

\paragraph{One endpoint.}
In the coefficient-one color, the edge is a natural or dual-natural
module of equal dimension, excluded by \Cref{w:natural-root}.
Otherwise the upper digit module, up to duality, is a symmetric power
of degree $c$ with $2\le c<p$. The lower natural or dual-natural
restriction has highest weight $\eta_0$, equal to $\eta$ or $\eta^*$,
and a highest simple-root string containing vectors $a_0,a_1,a_2$
of weights $\eta_0,\eta_0-\alpha,\eta_0-2\alpha$, since its restricted endpoint
coefficient is also greater than one. The distinct monomials
\[
 a_0^{c-1}a_2,\qquad a_0^{c-2}a_1^2
\]
both have weight $c\eta_0-2\alpha$ before the common twists. They
contradict the same single-root multiplicity bound. Thus every color
is impossible.
\end{proof}

\subsection{The signed series}

\begin{lemma}[Signed orbit cuts]\label{w:signed-cuts}
Let $J$ have sufficiently high type B, C, or D, and suppose $L_J(\eta)$ is self-dual, with $\eta\ne0$ restricted. In type D suppose for this lemma that no spin node of $\eta$ is active. Let a signed fundamental cut as in \Cref{w:universal-cuts} be given, including the last B and C nodes and excluding only the D-spin case, and let $U$ be its selected extremal orbit words. Use nonnegative dominant magnitudes, write $h$ for their nonzero count, and let $O^+$ be the unsigned permutation orbit.
\begin{enumerate}
\item If some unsigned word is unselected, then $U\subset O^+$. For $2\le h<e$, $|U|\le2$, with the binary possibility $\eta,s_h\eta$; for $h=e$, needed only in B and C, $U=\{\eta\}$.
\item If every unsigned word is selected and $h<e$, then either $h=1$ or $h=2$ with the two nonzero magnitudes equal.
\item If every unsigned word is selected and $h=e$ in B or C, all magnitudes are equal and $|U|\le3$.
\end{enumerate}
\end{lemma}
\begin{proof}
Now \(J\) is B,C or D of rank \(e\), and dominance in particular requires the prefix comparisons in the usual ordered signed coordinates. The orbit of \(\eta\) is antipodal by self-duality; make its dominant magnitudes nonnegative using the D diagram if necessary. First suppose no D spin nodes active (case of spin activity handled separately). For a universal fundamental cut above in \(\Omega\), upper orbit words \(U\) include \(\eta\), and each has all prefix sums nonnegative and is nonzero. Write \(h\) for the number of nonzero magnitudes in a word, and \(O^+\) for all unsigned permutations.

If some unsigned word is outside \(U\), comparison of total sums gives \(U\subset O^+\). If \(h<e\), we have all choices of signs (even in D, by zeros). Before the earliest negative word in decreasing lex order the first \(h-1\) positive coordinates are fixed, hence \(|U|\le e-h+1\). Indeed in this case all negative words are outside \(U\), and one can put the largest \(h-1\) magnitudes first positive, with the final nonzero one negative later; left prefix domination gives the lex bound. For \(2\le h<e\) in high rank this is at most half the unsigned words (there are at least \(\binom eh\)). In this range apply the A-word prefix description among those equal-total words. We obtain at most two upper words (the adj-shaped half for a word \(M,T^{e-2},m\) of three distinct magnitudes would be too large by \(e-h+1\), and the remaining binary singleton-extreme exception in this range has a singleton zero, already giving bound two by \(e-h+1\)). For a binary shape in that range with equal nonzero coordinates they are \(\eta\) and possibly \(s_h\eta\), also when \(h=e-1\) by left lex order. If instead \(h=e\), here only needed for B,C, the first \(e-1\) positives are fixed before the first negative word, giving just \(\eta\).

Suppose all unsigned words are selected. If $h<e$, every signed word
with positive first coordinate is forced in: otherwise it fails the
first-prefix comparison with a selected unsigned word starting with
zero. For $h\ge3$, put a smallest nonzero magnitude first positively
and two others negatively; the third prefix is negative. For $h=2$
with unequal magnitudes, the smaller positive followed by the larger
negative gives the same contradiction. Thus only $h=1$ or two equal
nonzero magnitudes can occur.

For $h=e$ in B or C, suppose the smallest and largest magnitudes are
$P<M$, and let $S$ be the sum of all $e$ magnitudes. The signed word
with $M$ first and all other entries negative must be selected, by
comparison with a selected unsigned word starting with $P$. The word
starting with $-P$ and having all other entries positive is unselected,
since its first prefix is negative. Domination of their total sums
would require
\[
 2M-S\ge S-2P,\qquad\text{hence}\qquad S\le M+P,
\]
impossible with at least three positive coordinates. All magnitudes
are therefore equal. A selected word cannot have two or more negatives,
since total-sum comparison would then force in every single-negative
word, including one with negative first prefix. Left lex order before
the first two-negative word now leaves only the all-positive word and
the words with one negative in either of the last two positions.
Thus $|U|\le3$.
\end{proof}

For a fixed series B, C, or D, use the following edge colors:
\begin{enumerate}
\item a single node $1$, with coefficient one or greater than one as separate colors;
\item a single interior nonspin node, separating index $2$ from larger indices;
\item the single last node in B or C;
\item any spin activity in D;
\item all remaining multiple support.
\end{enumerate}
Here ``interior'' in B and C excludes the last node. B and C are separate vertex types, and B's last-node color includes its spin weight. The D colors are invariant under interchange of its spin nodes.

\begin{proposition}[No homogeneous signed configuration]\label{w:signed-obstruction}
Four nodes of a protected run having one fixed series B, C, or D cannot have all pair edges of one color in the preceding palette.
\end{proposition}
\begin{proof}
Use the triple notation of \Cref{w:rational-composition}, and consider the possible common edge color.

\paragraph{D spin activity.} For D spin activity on the two edges, \(\eta\) has at least \(e-1\) nonzero coordinates. In its orbit take words starting with a smallest nonzero entry positive and two others negative, then many independent signs on the remaining nonzero ones respecting parity. If all $e$ coordinates are nonzero, parity leaves at least $2^{e-4}$ choices; if one is zero, it absorbs the parity adjustment and gives the same lower bound. Their first prefix is positive and third prefix negative. No two constructed words are antipodes, since the first coordinate is positive in each. Thus for $e\ge7$ this gives more than $e$ distinct antipodal pairs with mixed prefix signs. Self-duality supplies antipodality even in odd D rank, where it forces the last coordinate to be zero when opposition interchanges the spin nodes. Each pair appears (multiplicity one) as a natural ambient signed pair in \(\Omega\). In the spin maximum for \(\lambda\) as above, its chosen orientation therefore contributes a negative prefix somewhere. But at any one prefix at most one coordinate choice can have negative sum, by the pair-sum condition. Contradiction.

\paragraph{Remaining multiple support.} In the remaining multiple-support color, the orbit description forces \(U\) to have just \(\eta\). Indeed \(h=1\), two equal nonzero coordinates only, and all equal full magnitudes are all single-node shapes here. And a bound of two words in the description reduces to just \(\eta\) under multiple support: any additional one would have to dominate the active reflections outside, impossible (its difference down from highest must be supported only at each such node; if itself an active reflection use a different active node). Then every full selected weight also dominates all those reflections, giving just \(\eta\) with multiplicity one. This cannot accommodate the active indices of \(\lambda\).

\paragraph{One interior nonspin node.}
Write $\lambda=c\omega_a$ and $\eta=z\omega_h$, where $a$ and $h$
are interior nonspin nodes in their respective diagrams. Here
$2\le h<e$, with $h\le e-2$ in D. If all unsigned words were selected,
the two-equal-magnitude exception would force $h=2$. The palette then
also gives $a=2$, whereas the cut would contain at least
$\binom e2>2$ unsigned words. Thus only $\eta$, possibly together with
$s_h\eta$, occurs in the selected orbit.

As in type A, every full selected weight differs downward from $\eta$
only at $h$. If $s_h\eta$ is selected, use reflected words from two
distinct Dynkin neighbors; at the D fork either right neighbor suffices.
Intersecting their downward supports does not require the right bond
to be simple. The outside weight
$\eta-\alpha_h-\alpha_{h-1}$ is present by the left simple bond.
After the first lowering its vector is highest for that left rank-one
subgroup and has coefficient one there, so the second lowering is
nonzero in every characteristic. Domination of this outside weight
leaves at most the two entries $\eta,\eta-\alpha_h$, each of
multiplicity one. Since the cut has $a\ge2$ entries, it has exactly
two. The composed highest weight is a positive scalar multiple of
$2\eta-\alpha_h$, with both neighboring nodes active, contrary to
the direct edge's one-node digit shape.

\paragraph{The last B or C node.} For the last single node of B,C the fundamental maximum in the upper group orients every natural signed pair to the upper side. It therefore chooses half the orbit of \(\eta\), which has all equal positive magnitudes before sign changes. This contradicts the bound above.

\paragraph{Node one.} Coefficient one gives the prohibited natural embedding. For a single node 1 with both coefficients larger than one, write \(\lambda=z\omega_1\). Assign \(\epsilon_1,\epsilon_2,\epsilon_3\) at a maximum respectively to \(\eta,\eta-\alpha_1,\eta-2\alpha_1\), up to twist (distinct, nonzero and nonopposite, adjusting unused signs if needed). The top \(A_2\)-face in the upper module before twist is the natural symmetric power of degree \(z\). Thus the distinct weights \((z-1)\epsilon_1+\epsilon_3\) and \((z-2)\epsilon_1+2\epsilon_2\) both occur in these coordinates and restrict to $z\eta-2\alpha_1$ before the common twists. Their two weight spaces give multiplicity at least two along one simple-root ray from highest after restriction, contradicting \Cref{w:universal-cuts}.

All constructions require only fixed lower bounds on rank: four
coordinates suffice for the type A equal-gap test, three for the signed
total-sum comparison, and seven for the displayed D parity count.
They therefore apply after fixing one sufficiently large absolute
high-rank threshold.
\end{proof}

\subsection{The absolute path bound}

\begin{theorem}[Bound on protected high-type paths]\label{w:path-bound}
There is an absolute constant $D_0$ such that a prepared path all of whose socles are of high alternating or high classical type, and whose consecutive same-characteristic classical runs are protected, has fewer than $D_0$ nodes. This constant is independent of all chart dimensions, finite fields, and path-construction parameters.
\end{theorem}
\begin{proof}
Within a protected classical run, first color vertices by the four series A, B, C, and D. The edge palettes above have an absolute number of colors and are compatible with diagram changes. Finite Ramsey theory \citep{Ramsey1930} therefore bounds the number of nodes in the run: a sufficiently long run contains four same-series nodes whose pair edges are homogeneous, contrary to \Cref{w:a-obstruction,w:signed-obstruction}. All pairs of the selected nodes remain protected because the protection was imposed for every ancestor--descendant pair, not just consecutive nodes.

By \Cref{w:alternating-switch}, at most two high alternating nodes occur. By \Cref{w:characteristic-switch}, there cannot be two successive changes of characteristic among the classical nodes after ignoring the alternating interruptions. Thus there are only absolutely boundedly many same-characteristic classical runs. Applying the run bound to each gives $D_0$.
\end{proof}

\begin{proof}[Completion of the proof of \Cref{f:negative-family}]
Use the almost simple reduction of \Cref{f:thm-almost-simple-reduction} and the prepared path construction of \Cref{sec:patterns}. Choose the high thresholds and $D_0$ first, then take the number $D$ of reserved generations greater than $D_0$, with room for the initial displacement and one remaining bounded-type stage. Choose the lattice field size $q$ above the finitely many absolute line requirements. Next choose the terminal chart increment $r$ using \Cref{w:selection} and \Cref{w:bounded-exclusion}. Choose the earlier trial increments backwards from $r$ as in \Cref{sec:patterns}, then $m$ above their ranks and $n$ large enough for all residual interval lengths and bounded spans. Finally choose $N$ for the initial pattern demand and enough private charts $g$ for all finite losses.

At each stage, \Cref{sec:classical,sec:alternating} supplies a pruned chart with the ancestor properties of \Cref{w:prepared-path}. A bounded-order socle is excluded by the residual interval length. A bounded-type Lie socle is impossible by \Cref{w:bounded-exclusion}. In every remaining case the next high-type node can be selected, with \Cref{w:selection} adding the protections whenever the current classical run continues. This produces more than $D_0$ high-type nodes, contradicting \Cref{w:path-bound}. It proves the required finite large-family nonrepresentation statement. Together with \Cref{sec:small}, this gives a finite minimum strictly greater than seven. Its explicit size ceiling and the description of $L_0$ are supplied in \Cref{sec:quantitative,sec:minimum} after the uniform-decision argument.
\end{proof}
